\documentclass[10pt,a4paper]{amsart}
\usepackage{amsmath,amssymb,amsthm,graphicx}
\usepackage{verbatim}
\newenvironment{reference}{\comment}{\endcomment}
%\newenvironment{reference}{}{}
\newenvironment{slogan}{\comment}{\endcomment}
\newenvironment{history}{\comment}{\endcomment}
% For commutative diagrams we use Xy-pic
\usepackage[all]{xy}

% We use 2cell for 2-commutative diagrams.
\xyoption{2cell}
% Theorem environments.
%
\theoremstyle{plain}
\newtheorem{theorem}[subsection]{Theorem}
\newtheorem{proposition}[subsection]{Proposition}
\newtheorem{lemma}[subsection]{Lemma}

\theoremstyle{definition}
\newtheorem{definition}[subsection]{Definition}
\newtheorem{example}[subsection]{Example}
\newtheorem{exercise}[subsection]{Exercise}
\newtheorem{situation}[subsection]{Situation}

\theoremstyle{remark}
\newtheorem{remark}[subsection]{Remark}
\newtheorem{remarks}[subsection]{Remarks}

\numberwithin{equation}{subsection}

% Macros
%
\def\lim{\mathop{\mathrm{lim}}\nolimits}
\def\colim{\mathop{\mathrm{colim}}\nolimits}
\def\Spec{\mathop{\mathrm{Spec}}}
\def\Hom{\mathop{\mathrm{Hom}}\nolimits}
\def\Ext{\mathop{\mathrm{Ext}}\nolimits}
\def\SheafHom{\mathop{\mathcal{H}\!\mathit{om}}\nolimits}
\def\SheafExt{\mathop{\mathcal{E}\!\mathit{xt}}\nolimits}
\def\Sch{\mathit{Sch}}
\def\Mor{\mathop{\mathrm{Mor}}\nolimits}
\def\Ob{\mathop{\mathrm{Ob}}\nolimits}
\def\Sh{\mathop{\mathit{Sh}}\nolimits}
\def\NL{\mathop{N\!L}\nolimits}
\def\CH{\mathop{\mathrm{CH}}\nolimits}
\def\proetale{{pro\text{-}\acute{e}tale}}
\def\etale{{\acute{e}tale}}
\def\QCoh{\mathit{QCoh}}
\def\Ker{\mathop{\mathrm{Ker}}}
\def\Im{\mathop{\mathrm{Im}}}
\def\Coker{\mathop{\mathrm{Coker}}}
\def\Coim{\mathop{\mathrm{Coim}}}

% Boxtimes
%
\DeclareMathSymbol{\boxtimes}{\mathbin}{AMSa}{"02}

%
% Macros for moduli stacks/spaces
%
\def\QCohstack{\mathcal{QC}\!\mathit{oh}}
\def\Cohstack{\mathcal{C}\!\mathit{oh}}
\def\Spacesstack{\mathcal{S}\!\mathit{paces}}
\def\Quotfunctor{\mathrm{Quot}}
\def\Hilbfunctor{\mathrm{Hilb}}
\def\Curvesstack{\mathcal{C}\!\mathit{urves}}
\def\Polarizedstack{\mathcal{P}\!\mathit{olarized}}
\def\Complexesstack{\mathcal{C}\!\mathit{omplexes}}
% \Pic is the operator that assigns to X its picard group, usage \Pic(X)
% \Picardstack_{X/B} denotes the Picard stack of X over B
% \Picardfunctor_{X/B} denotes the Picard functor of X over B
\def\Pic{\mathop{\mathrm{Pic}}\nolimits}
\def\Picardstack{\mathcal{P}\!\mathit{ic}}
\def\Picardfunctor{\mathrm{Pic}}
\def\Deformationcategory{\mathcal{D}\!\mathit{ef}}
\usepackage{cleveref}
\crefname{originalsection}{original section}{original sections}
\Crefname{originalsection}{Original section}{Original sections}
\makeatletter
\newlabel{example-stone-Cech-discrete}{{example-stone-Cech-discrete}{source}{}{sourcepointer-1}{}}
\newlabel{example-stone-Cech-discrete@cref}{{[originalsection][0][]example-stone-Cech-discrete}{[1][0][]source}{}{}{}}
\makeatother
\crefname{originallemma}{original lemma}{original lemmas}
\Crefname{originallemma}{Original lemma}{Original lemmas}
\makeatletter
\newlabel{lemma-closed-map}{{lemma-closed-map}{source}{}{sourcepointer-2}{}}
\newlabel{lemma-closed-map@cref}{{[originallemma][0][]lemma-closed-map}{[1][0][]source}{}{}{}}
\makeatother
\crefname{originallemma}{original lemma}{original lemmas}
\Crefname{originallemma}{Original lemma}{Original lemmas}
\makeatletter
\newlabel{lemma-find-compact-subset}{{lemma-find-compact-subset}{source}{}{sourcepointer-3}{}}
\newlabel{lemma-find-compact-subset@cref}{{[originallemma][0][]lemma-find-compact-subset}{[1][0][]source}{}{}{}}
\makeatother
\crefname{originallemma}{original lemma}{original lemmas}
\Crefname{originallemma}{Original lemma}{Original lemmas}
\makeatletter
\newlabel{lemma-isomorphism}{{lemma-isomorphism}{source}{}{sourcepointer-4}{}}
\newlabel{lemma-isomorphism@cref}{{[originallemma][0][]lemma-isomorphism}{[1][0][]source}{}{}{}}
\makeatother
\crefname{originallemma}{original lemma}{original lemmas}
\Crefname{originallemma}{Original lemma}{Original lemmas}
\makeatletter
\newlabel{lemma-rainwater}{{lemma-rainwater}{source}{}{sourcepointer-5}{}}
\newlabel{lemma-rainwater@cref}{{[originallemma][0][]lemma-rainwater}{[1][0][]source}{}{}{}}
\makeatother
\crefname{originalproposition}{original proposition}{original propositions}
\Crefname{originalproposition}{Original proposition}{Original propositions}
\makeatletter
\newlabel{proposition-projective-in-category-hausdorff-qc}{{proposition-projective-in-category-hausdorff-qc}{source}{}{sourcepointer-6}{}}
\newlabel{proposition-projective-in-category-hausdorff-qc@cref}{{[originalproposition][0][]proposition-projective-in-category-hausdorff-qc}{[1][0][]source}{}{}{}}
\makeatother
\title{Every quasi-compact Hausdorff space has a canonical minimal extremally disconnected cover, unique up to isomorphism}
\author{The Stacks Project Authors}
\date{stacks-project:topology.tex:local-snapshot-20261012; GNU FDL 1.2 or later}
\begin{document}
\maketitle
\noindent\textit{Editorial note.} All delivered mathematical text is the author's own, copied from the pinned Stacks Project chapter source sources/stacks-topology/pinned-original-source.tex; no statement or proof step was rewritten, reordered or removed.\par\medskip
\noindent\textit{Editorial note.} Preamble provenance: the Stacks Project keeps one single shared preamble for every chapter; this card takes it directly from the staged sources/stacks-topology/preamble.tex as three verbatim spans (lines 8-12, 14-18 and 154-217), so no derived file was written and nothing was added or redefined. The wrapper class and the base packages are supplied by the assembler.\par\medskip
\noindent\textit{Editorial note.} Bibliography provenance: neither the counted statement nor its proof contains a citation (this lane rescans the whole span before writing the card), so the card carries no bibliography and no entry was reconstructed.\par\medskip
\noindent\textit{Editorial note.} Selection provenance: this card was opened under the file-based control contract (CONTROL.md), with the mechanical pre-card gate run on its statement and proof interval before the card was written.\par\medskip
\noindent\textit{Allowed standard prerequisites (editorial).}\begin{itemize}
\item The description of the Stone-Cech compactification of a discrete space as extremally disconnected. Invoked inside the counted proof (source lines 5627-5632). Produces the quasi-compact Hausdorff extremally disconnected space covering the target.
\item The existence of a minimal quasi-compact subset surjecting onto the target. Invoked inside the counted proof (source lines 5633-5635). Gives the minimal cover used for uniqueness.
\item The lifting property of extremally disconnected quasi-compact Hausdorff spaces, and the statement that a surjection over an extremally disconnected target which is minimal is a homeomorphism. Invoked inside the counted proof (source lines 5636-5645 and 5650-5658). Produces the self-map and the isomorphism in the uniqueness argument.
\item The statement that a non-identity surjective self-map of a Hausdorff space admits a proper closed subset covering the space with its image. Invoked inside the counted proof (source lines 5641-5643). Rules out a non-identity self-map of the minimal cover.
\item The statement that a continuous map from a quasi-compact to a Hausdorff space is closed. Invoked inside the counted proof (source lines 5651-5652). Makes the comparison map closed.
\end{itemize}\medskip
\ifdefined\hypertarget\hypertarget{sourcepointer-1}{}\fi
\noindent\textit{Source reference.} Original section example-stone-Cech-discrete: sources/stacks-topology/pinned-original-source.tex, source lines 5590--5590. Navigation pointer only: the counted text refers to the Stacks item example-stone-Cech-discrete, which is an existing labeled result of this same chapter and is not part of this card. The pointer carries the target's own label as its number and its definition line as its span; its content is not reproduced here.\par\medskip
\ifdefined\hypertarget\hypertarget{sourcepointer-2}{}\fi
\noindent\textit{Source reference.} Original lemma lemma-closed-map: sources/stacks-topology/pinned-original-source.tex, source lines 3474--3474. Navigation pointer only: the counted text refers to the Stacks item lemma-closed-map, which is an existing labeled result of this same chapter and is not part of this card. The pointer carries the target's own label as its number and its definition line as its span; its content is not reproduced here.\par\medskip
\ifdefined\hypertarget\hypertarget{sourcepointer-3}{}\fi
\noindent\textit{Source reference.} Original lemma lemma-find-compact-subset: sources/stacks-topology/pinned-original-source.tex, source lines 5503--5503. Navigation pointer only: the counted text refers to the Stacks item lemma-find-compact-subset, which is an existing labeled result of this same chapter and is not part of this card. The pointer carries the target's own label as its number and its definition line as its span; its content is not reproduced here.\par\medskip
\ifdefined\hypertarget\hypertarget{sourcepointer-4}{}\fi
\noindent\textit{Source reference.} Original lemma lemma-isomorphism: sources/stacks-topology/pinned-original-source.tex, source lines 5481--5481. Navigation pointer only: the counted text refers to the Stacks item lemma-isomorphism, which is an existing labeled result of this same chapter and is not part of this card. The pointer carries the target's own label as its number and its definition line as its span; its content is not reproduced here.\par\medskip
\ifdefined\hypertarget\hypertarget{sourcepointer-5}{}\fi
\noindent\textit{Source reference.} Original lemma lemma-rainwater: sources/stacks-topology/pinned-original-source.tex, source lines 5573--5573. Navigation pointer only: the counted text refers to the Stacks item lemma-rainwater, which is an existing labeled result of this same chapter and is not part of this card. The pointer carries the target's own label as its number and its definition line as its span; its content is not reproduced here.\par\medskip
\ifdefined\hypertarget\hypertarget{sourcepointer-6}{}\fi
\noindent\textit{Source reference.} Original proposition proposition-projective-in-category-hausdorff-qc: sources/stacks-topology/pinned-original-source.tex, source lines 5530--5530. Navigation pointer only: the counted text refers to the Stacks item proposition-projective-in-category-hausdorff-qc, which is an existing labeled result of this same chapter and is not part of this card. The pointer carries the target's own label as its number and its definition line as its span; its content is not reproduced here.\par\medskip
\begin{lemma}
\label{lemma-existence-projective-cover}
\begin{slogan}
Every quasi-compact Hausdorff space has a canonical
extremally disconnected cover
\end{slogan}
Let $X$ be a quasi-compact Hausdorff space.
There exists a continuous surjection $X' \to X$ with $X'$
quasi-compact, Hausdorff, and extremally disconnected.
If we require that every proper closed subset of $X'$ does not
map onto $X$, then $X'$ is unique up to isomorphism.
\end{lemma}

\begin{proof}
Let $Y = X$ but endowed with the discrete topology. Let $X' = \beta(Y)$.
The continuous map $Y \to X$ factors as $Y \to X' \to X$. This
proves the first statement of the lemma by
Example \ref{example-stone-Cech-discrete}.

\medskip\noindent
By Lemma \ref{lemma-find-compact-subset} we can find a quasi-compact subset
$E \subset X'$ surjecting onto $X$ such that no proper closed subset of $E$
surjects onto $X$.
Because $X'$ is extremally disconnected there exists a continuous map
$f : X' \to E$ over $X$
(Proposition \ref{proposition-projective-in-category-hausdorff-qc}).
Composing $f$ with the map $E \to X'$ gives a continuous selfmap
$f|_E : E \to E$. Observe that $f|_E$ has to be surjective as otherwise
the image would be a proper closed subset surjecting onto $X$.
Hence $f|_E$ has to be $\text{id}_E$ as otherwise
Lemma \ref{lemma-rainwater} shows that $E$ isn't minimal.
Thus the $\text{id}_E$ factors through the extremally disconnected
space $X'$. A formal, categorical argument (using the characterization of
Proposition \ref{proposition-projective-in-category-hausdorff-qc})
shows that $E$ is extremally disconnected.

\medskip\noindent
To prove uniqueness, suppose we have a second $X'' \to X$
minimal cover. By the lifting property proven in
Proposition \ref{proposition-projective-in-category-hausdorff-qc}
we can find a continuous map $g : X' \to X''$ over $X$.
Observe that $g$ is a closed map (Lemma \ref{lemma-closed-map}).
Hence $g(X') \subset X''$ is a closed subset surjecting onto $X$
and we conclude $g(X') = X''$ by minimality of $X''$.
On the other hand, if $E \subset X'$ is a proper closed subset,
then $g(E) \not = X''$ as $E$ does not map onto $X$ by minimality
of $X'$. By Lemma \ref{lemma-isomorphism} we see that $g$ is an isomorphism.
\end{proof}

\end{document}
